% Generated draft: requires mathematical and source review.
\documentclass[11pt]{article}
\usepackage[margin=25mm]{geometry}
\usepackage{fontspec}
\setmainfont{Latin Modern Roman}
\usepackage{amsmath,amssymb,mathtools,unicode-math}
\setmathfont{Latin Modern Math}
\usepackage{enumitem,needspace,xurl,hyperref}
\setlength{\parindent}{0pt}
\setlength{\parskip}{5pt}
\setlength{\emergencystretch}{3em}
\newcommand{\N}{\mathbb{N}}
\newcommand{\Z}{\mathbb{Z}}
\newcommand{\Q}{\mathbb{Q}}
\newcommand{\R}{\mathbb{R}}
\newcommand{\C}{\mathbb{C}}
\newcommand{\F}{\mathbb{F}}
\newcommand{\A}{\mathbb{A}}
\newcommand{\PP}{\mathbb{P}}
\newcommand{\E}{\mathbb{E}}
\newcommand{\eps}{\varepsilon}
\newcommand{\dd}{\,\mathrm{d}}
\begin{document}
\section*{UnsolvedMath ID: 20003350}
\textbf{Source problem code:} catalogue entry
\subsection{Definitions and mathematical statement}
In the standard pure-Shimura interpretation, let $S$ be a Shimura variety over $\C$. A special subvariety is an irreducible component arising from a Shimura subdatum and the permitted Hecke translates; special points are the zero-dimensional cases, equivalently CM points in this setting. Conjecture 51 of Y. Andr\'e and F. Oort asserts that for every subset $\Sigma$ of special points,
\[
 \overline\Sigma^{\mathrm{Zar}}\text{ is a finite union of special subvarieties}.
\]
Equivalently, every irreducible subvariety with Zariski-dense special points is special. The assigned extraction suppresses the ambient variety and ends mid-sentence in its literature discussion; no mixed-Shimura extension is intended here.

The supplied status assessment distinguishes three levels. It reports that the missing historical continuation says the one-dimensional-closure case was known in 2004 under the intended Generalized Riemann Hypothesis. Klingler--Yafaev's published conditional theorem covers arbitrary pure Shimura varieties under the GRH hypotheses of that theorem. Tsimerman's published result proves the assertion unconditionally for $\mathcal A_g$, with standard Hodge-type and abelian-type consequences. Pila--Shankar--Tsimerman, \texttt{arXiv:2109.08788v4}, states the unconditional maximal-special-subvariety theorem for arbitrary Shimura varieties, but the literature located in the supplied check still treated it as a preprint. These publication and conditionality qualifications accompany the registry solved label and are not replaced by an independently updated status claim.

The attempt additionally gives an explicit elementary family in $Y(1)^n\cong\A^n_{\C}$, where $Y(1)$ is the elliptic $j$-line without its cusp and singular moduli are its CM $j$-values. Let
\[
 \Sigma=\bigcup_{\ell=1}^{r}\prod_{i=1}^n A_{\ell i},
\]
where each $A_{\ell i}$ is a set of singular moduli and $r$ is finite. Omit any rectangle with an empty factor. For each remaining rectangle let $F_\ell=\{i:A_{\ell i}\text{ is finite}\}$, and for each choice $a\in\prod_{i\in F_\ell}A_{\ell i}$ put
\[
 X(F_\ell,a)=\{z\in Y(1)^n:z_i=a_i\text{ for }i\in F_\ell\}.
\]
The aid's exact closure formula is
\[
 \overline\Sigma^{\mathrm{Zar}}=
       \bigcup_{\ell}\ \bigcup_{a\in\prod_{i\in F_\ell}A_{\ell i}}
                 X(F_\ell,a).
\]
An infinite subset of $Y(1)$ is Zariski dense, while a finite one remains its singleton choices, explaining the product closure. Each slice is special and irreducible, of dimension $n-|F_\ell|$. The irreducible components are exactly the distinct inclusion-maximal slices, using
\[
 X(F,a)\subset X(F',a')
 \quad\Longleftrightarrow\quad F'\subset F\text{ and }a_i=a'_i\ (i\in F').
\]
For nonempty $\Sigma$, the closure dimension is $\max_\ell(n-|F_\ell|)$; for an empty union the closure and component set are empty. This classifies finite unions of Cartesian CM families only, not correlated CM sets or the general Andr\'e--Oort theorem.
\subsection{Short English statement}
Andr\'e--Oort asks for special irreducible components in the Zariski closure of any set of special points. The supplied status distinguishes conditional and published results from a general preprint theorem. Its extra worked case computes closures of finite unions of Cartesian CM families in powers of the $j$-line as explicit coordinate special slices.
\subsection{Sources}
\begin{itemize}
\item Local UnsolvedMath registry, permanent ID $20003350$, source code catalogue entry; canonical location \texttt{aim-other-notes.json}, \texttt{notes[156]}.
\item AIM \emph{Emerging applications of measure rigidity}, Conjecture 51 (Y. Andr\'e--F. Oort [3,153]), \url{https://aimath.org/WWN/measrigid/measrigid.pdf}. It cites [45,46] for historical partial results; full entries are absent. The supplied aid recovers the truncated one-dimensional GRH continuation.
\item Klingler--Yafaev's published GRH-conditional theorem, Tsimerman's published $\mathcal A_g$ theorem, and Pila--Shankar--Tsimerman \texttt{arXiv:2109.08788v4}, as described in the attached summary. The supplied literature check labels the last item preprint-only; no additional publication details are invented.
\item Research attempt 1 is marked partial result, proved, and low novelty confidence, despite the registry solved-in-literature label, with an explicit machine-generated, non-peer-reviewed research-aid caveat requiring independent verification before citation.
\end{itemize}
\end{document}
